% Introduction — draft by M3 (2026-08-03), empirical claims rewritten by M1 (2026-08-05)
% after the natural-speech scale results came in. See docs/decisions.md for what changed and why.

\section{Introduction}
\label{sec:introduction}

Large Audio-Language Models (LALMs), or Speech LLMs \cite{chu2024qwen2audio, radford2022whisper}, have rapidly advanced the paradigm of voice-based artificial intelligence. By integrating neural audio encoders directly into frozen or instruction-tuned text LLM backbones, these end-to-end architectures eliminate the cascade latency of traditional Automatic Speech Recognition (ASR) pipelines. Consequently, modern Speech LLMs demonstrate remarkable capabilities across spoken question answering, dialogue generation, and audio reasoning \cite{sakshi2024mmau, wang2025mmsu}. However, despite their surface-level fluency, these models implicitly operate under a brittle closed-world assumption: they assume that every user query is fully grounded in the provided acoustic context and that a factual answer always exists.

In real-world deployment, acoustic evidence is frequently noisy, incomplete, or entirely absent regarding a specific question. When faced with such unanswerable queries, contemporary Speech LLMs exhibit severe unreliability. Rather than expressing uncertainty or abstaining, models default to an affirmative forced-choice bias, producing fluent yet unfounded \textit{hallucinations} \cite{zhao2026halluaudio, kuan2026aquabench}. Naïve alignment interventions---such as prompting the model with ``I don't know'' (IDK) instructions---often fail to solve this issue; instead of fostering true epistemic awareness, they trigger indiscriminate over-refusal, causing models to reject valid questions where sufficient evidence is clearly present \cite{ma2025reliable}.

This failure raises a fundamental diagnostic question: \textbf{do Speech LLMs hallucinate because they fail to perceive the acoustic signal (``a hearing failure''), or because they fail to reason about what the evidence establishes (``an epistemic reasoning failure'')?} To answer it, we compare an end-to-end Speech LLM against an ASR-cascade baseline (Whisper $\rightarrow$ text LLM) operating on the \textit{exact same audio}, across a structured taxonomy of answerable (stated, inferable) and unanswerable spoken queries.

On our first diagnostic set---100 hand-verified items over synthesized (TTS) speech---the answer appeared unambiguous. Qwen2-Audio hallucinated on 92.5\% of unanswerable questions under standard prompting; IDK-prompting cut this to 17.5\%, but at the cost of a 61.7\% over-refusal rate on answerable ones. The cascade, reading a Whisper transcript of the same recordings, reached 2.5\% hallucination at only 6.7\% over-refusal. Since the cascade discriminates evidence almost perfectly from transcribed text alone, the necessary information is demonstrably recoverable from the audio, and the bottleneck would seem to lie in epistemic reasoning rather than perception.

\textbf{This conclusion does not survive contact with natural speech.} We re-ran the identical protocol on 376 native SQuAD~2.0 questions---human-written, adversarially designed to look answerable---grounded in recordings read aloud by 40 human speakers (NMSQA). Here the cascade's advantage disappears entirely: it hallucinates on 87.2\% of unanswerable questions under standard prompting and still on 67.7\% under IDK-prompting, statistically indistinguishable from the end-to-end model's 69.8\% and now at a higher over-refusal cost (11.3\% versus 0.7\%). Two conclusions follow, both cautionary. First, the apparent superiority of the cascade was largely an artifact of synthesized speech: TTS audio is cleaner and more predictable than human speech, and a diagnostic run on it overstates how much an ASR front-end recovers. Second, IDK-prompting does not scale from generated to native adversarial questions---what removed three quarters of the hallucinations on our pilot removes only a quarter of them here. The perception-versus-reasoning dichotomy is therefore not resolved by the cascade comparison; both pipelines fail on the same items, and the failure is systemic rather than attributable to either component in isolation.

This makes lightweight, prompt-level fixes an unpromising direction and motivates looking inside the model instead. Existing efforts to address LALM unreliability either focus solely on diagnostic measurement \cite{kuan2026aquabench, zhao2026halluaudio} or rely on post-generation sampling and post-hoc self-verification \cite{kuan2026uncertainty}. While post-generation methods (e.g., semantic entropy over $K$ sampled outputs) can detect uncertainty, they introduce substantial computation cost and decoding latency that undermine the real-time advantages of voice interfaces. In this work we ask whether a Speech LLM's internal representation state contains sufficient signal to detect unanswerability \textit{before} token generation begins. Specifically, we transfer pre-generation attention probing \cite{miftakhova2026probing} to the multimodal setting: a lightweight probe is trained on frozen prompt representations---leaving the LALM's own weights untouched---and used as an early-stage abstention gate.

Our primary contributions are three-fold:
\begin{enumerate}
    \item \textbf{A diagnostic set and a natural-speech scale validation.} We introduce a fine-grained spoken QA diagnostic set (Pilot-100) with structured unanswerability subtypes (absent entities, missing attributes, false premises, off-topic queries), and validate all findings at scale on 376 native SQuAD~2.0 questions grounded in human-read speech, with no LLM-generated questions in the validation set.
    \item \textbf{Evidence that TTS-based diagnosis is misleading.} We show that conclusions drawn from synthesized speech do not transfer to human speech: an ASR-cascade that appears to solve unanswerability on TTS audio (2.5\% hallucination) fails on natural recordings (67.7\%), and prompt-level abstention instructions lose most of their effect once questions are human-written rather than generated. This reframes a result that the pilot alone would have reported as a perception-versus-reasoning dissociation.
    \item \textbf{Pre-generation probing as a mitigation.} We propose a mitigation that leaves the model's weights unchanged: a probe over pre-generation hidden states detects unanswerability before decoding and gates abstention, avoiding the sampling cost of post-hoc uncertainty estimation.
\end{enumerate}
